\section*{Acknowledgements}

% Not compulsory. Acknowledgements should be brief, and should not include thanks to anonymous referees and editors, or effusive comments. Grant or contribution numbers may be acknowledged.

% \TODO{Following seems effusive, should pare it down}
% The NeuroBench project represents a collaborative and inclusive effort, wherein contributions from a diverse group of researchers from both academia and industry have been instrumental in shaping the benchmark design and motivating its release. The authors listed in this paper are only a small subset of the larger community whose collective expertise and dedication have been integral to the development of the NeuroBench initiative. Our project also builds off prior efforts and collaborations from Telluride, Capo Caccia, and NICE, which have driven the NeuroBench community development and trailblazed paths towards inclusive benchmark design. We are immensely grateful to the NeuroBench community for their wisdom, effort, and passion, which they have shared with us in the form of countless meetings, research publications, and other forms of engagement. Without their support and input, NeuroBench would not have been matured from concept to a full-scale community-driven effort, and we look forward to continued collaboration with the community in the evolution of the project.

This work has been supported in parts by Semiconductor Research Corporation (SRC), the European Research Council (ERC) under the European Union’s Horizon 2020 research and innovation programme (grant agreement No.~101001448), a grant from the Research Grants Council of the Hong Kong Special Administrative Region, China [Project No.~CityU 11200922], ARC Laureate Fellowship FL210100156,  and the EU H2020 project BeFerroSynaptic (871737). We acknowledge the financial support of the CogniGron research center and the Ubbo Emmius Funds (Univ. of Groningen). We acknowledge a contribution from the Italian National Recovery and Resilience Plan (NRRP), M4C2, funded by the European Union –NextGenerationEU (Project IR0000011, CUP B51E22000150006, “EBRAINS-Italy”). The work of SynSense was partially supported by the European Commission, under the Horizon grant Ferro4Edge AI (grant agreement 101135656).This work is partly funded by the German Federal Ministry of Education and Research (BMBF) and the free state of Saxony within the ScaDS.AI center of excellence for AI research and by the German Federal Ministry for Economic Affairs and Climate Action (BMWK) under contract 01MN23004F (ESCADE).

Sandia National Laboratories is a multi-mission laboratory managed and operated by National Technology \& Engineering Solutions of Sandia, LLC (NTESS), a wholly owned subsidiary of Honeywell International Inc., for the U.S. Department of Energy’s National Nuclear Security Administration (DOE/NNSA) under contract DE-NA0003525. This written work is authored by an employee of NTESS. The employee, not NTESS, owns the right, title and interest in and to the written work and is responsible for its contents. Any subjective views or opinions that might be expressed in the written work do not necessarily represent the views of the U.S. Government. The publisher acknowledges that the U.S. Government retains a non-exclusive, paid-up, irrevocable, world-wide license to publish or reproduce the published form of this written work or allow others to do so, for U.S. Government purposes. The DOE will provide public access to results of federally sponsored research in accordance with the DOE Public Access Plan. This paper describes objective technical results and analysis. Any subjective views or opinions that might be expressed in the paper do not necessarily represent the views of the U.S. Department of Energy or the United States Government.

\section*{Author contributions statement}
Authors are grouped based on contributions, and ordered alphabetically within groups.
JY led project discussions and management. JY and KVdB implemented harness metric infrastructure, conducted experiments, and prepared the manuscript.
The following authors primarily developed the main algorithm track results: DdB and MF on the keyword few-shot class-incremental learning task; GT and SW on the event camera object detection task; PH, PVS and BZ on the non-human primate motor prediction task; YB and DK on the chaotic function prediction task; and NP led harness infrastructure development. The following authors primarily developed the main system track results: WK and MAK on the acoustic scene classification task; AP and PS on the QUBO task.
SHA, GVJ, BL, AM, AKM, GL, and TS developed components of the algorithm track harness infrastructure.
ZA, MA, BA, AGA, CB, AB, PB, SB, SB, GC, EC, FC, GdC, AD, AD, MD, YD, JE, TF, JF, VF, SF, PMF, WG, AG, HAG, GI, SJ, VK, LK, JCK, LK, RK, DK, YL, SL, HM, RM, FMM, CM, KM, DRM, EN, TN, FO, AO, PP, JP, MP, CP, MAP, CP, AR, YS, CJSS, AvS, JS, SS, CS, JS, SS, SBS, MS, AS, MS, KS, TCS, JT, NT, GU, MV, CMV, BV, AY, and FTZ participated in discussions during meetings, prepared sections for the present manuscript and/or its preprint, and reviewed the manuscript.
CF and VJR jointly supervised the project, reviewed the manuscript, and analyzed results.

% Must include all authors, identified by initials, for example:
% A.A. conceived the experiment(s),  A.A. and B.A. conducted the experiment(s), C.A. and D.A. analysed the results.  All authors reviewed the manuscript. 